\chapter{Related Work}
\label{sec:related}

\noindent \textbf{Parametric optimization:} For some flash loan attack cases,
researchers~\cite{cao2021flashot, qin2021attacking} manually extracted formulas
of state transition functions of the DeFi smart contracts, defined related
parameter constraints, and used an off-the-shelf optimizer to search for
parameters which yield the best profit. However, this technique requires
significant manual efforts and expert knowledge of the underlying DeFi
protocols. Furthermore, unlike {\name}, this technique also requires prior
knowledge of the sequence of attack actions. It is therefore not suitable
%is sufficient to check whether an attack vector is optimal but might not be
%enough 
to detect the possibility of flash attacks before the attacks happen.

\noindent \textbf{Static Analysis:} 
Slither~\cite{feist2019slither}, Securify~\cite{tsankov2018securify},
Zeus~\cite{kalra2018zeus}, Park~\cite{zheng2022park} and
SmartCheck~\cite{tikhomirov2018smartcheck} apply static analysis techniques to
verify smart contracts. 
%They detect pre-defined patterns of vulnerabilities and
%over-approximate program states, which inevitably cause false positives and
%negatives. 
There are also several works that use
symbolic execution~\cite{king1976symbolic} to explore the program states of a
smart contract, looking for an execution path that violates a user-defined
invariant, e.g., Mythril~\cite{mythril}, Oyente~\cite{luu2016making},
FairCon~\cite{liu2020towards}, ETHBMC~\cite{frank2020ethbmc},
SmartCopy~\cite{feng2019precise}, and Manticore~\cite{mossberg2019manticore}.
These techniques tend to operate with one contract at a time and therefore cannot handle flash loan attacks that involve multiple contracts.
When analyzing traces of DeFi contracts, these techniques also may suffer from the
path explosion and the complicated logics of the DeFi contracts.
{\name} uses its novel synthesis-via-approximation techniques to avoid these issues.

%These techniques suffer from the path explosion and the complicated logics 
%and cannot scale to 
%of DeFi contracts.
%Also both 
%
%However, as we showed in Section~\ref{sec:example}, symbolic execution tools do
%not scale well and suffer from path explosion. 
%Slither~\cite{feist2019slither}, Securify~\cite{tsankov2018securify},
%Zeus~\cite{kalra2018zeus}, Park~\cite{zheng2022park} and
%SmartCheck~\cite{tikhomirov2018smartcheck} apply static analysis techniques to
%verify smart contracts. They detect pre-defined patterns of vulnerabilities and
%over-approximate program states, which inevitably cause false positives and
%negatives. They are suitable for local analysis of contracts, but not suitable
%for analysis that must consider interactions between contracts and on-chain
%states. 

\noindent \textbf{Fuzzing smart contracts:}
ContractFuzzer~\cite{jiang2018contractfuzzer}, sFuzz~\cite{nguyen2020sfuzz},
ContraMaster~\cite{wang2020oracle} and SMARTIAN~\cite{choi2021smartian}
introduce new fuzzing techniques to discover vulnerabilities in smart
contracts. However, these techniques either only work on one contract or focus
on specific vulnerabilities like re-entrancy. Almost all flash loan attacks
interact with multiple contracts with complicated logics. Moreover, for many
flash loan attacks the action parameters have to be carefully chosen to
generate positive profits. Such a parameter combination is unlikely to be found
by random fuzzing.
